\documentclass[9pt]{IEEEtran}

\usepackage[english]{babel}
\usepackage{graphicx}
\usepackage{url}
\usepackage{booktabs}
\usepackage{amsmath}
\usepackage{hyperref}
\usepackage{float}
\usepackage[T1]{fontenc}
\usepackage[utf8]{inputenc}
\usepackage{lmodern}
\input{glyphtounicode}
\pdfgentounicode=1

\title{\vspace{0ex}
Programming Assignment 1: Preferential Web Crawler for MedOverNet}

\author{Tibor Koderman, Jerneja Krajcar, Dino Džaferagić}\vspace{-4.0ex}

\begin{document}
\maketitle

\section{Introduction}
In this project, we built a standalone preferential crawler for the MedOverNet domain (\url{https://medover.zurnal24.si/}). The goal is to crawl relevant pages first, respect robots rules, avoid duplicate content, and store crawl data in PostgreSQL using the required crawldb schema.

The crawler is designed as a multi-worker system with a database-backed frontier. Each worker takes one URL, checks robots and politeness limits, downloads or renders the page, detects duplicates, stores the result, extracts new links and images, scores discovered URLs, and pushes them back to the frontier.

This report is a first draft foundation. We describe implemented components and design decisions. Where final run outputs are still missing, we mark clear placeholders for later completion.

\section{System Architecture and Implementation}
\subsection{High-level pipeline}
The crawler processes each URL in the following order:
\begin{enumerate}
    \item Claim next URL from the frontier (highest score first).
    \item Skip re-download if the URL was already processed (only update link metadata).
    \item Check robots rules and enforce politeness delay.
    \item Download content (or render with headless browser when needed).
    \item Detect exact duplicates using content hash.
    \item Store page data and metadata in crawldb.
    \item Parse outgoing links and images.
    \item Score discovered URLs and enqueue them into frontier.
\end{enumerate}

\subsection{Main components}
The implementation is split into clear modules:
\begin{itemize}
    \item Worker orchestration and crawl loop: \texttt{scheduler.py}, \texttt{run\_crawler.py}.
    \item Frontier and page persistence: \texttt{frontier\_store.py}, \texttt{page\_store.py}.
    \item Downloader and renderer: \texttt{downloader.py}, \texttt{renderer.py}.
    \item Parsing links and images: \texttt{parser.py}.
    \item Canonicalization and dedup utils: \texttt{url\_canonicalizer.py}, \texttt{dedup.py}.
    \item Robots and politeness: \texttt{robots.py}, \texttt{politeness.py}.
\end{itemize}

\subsection{Division of work}
\textbf{Implemented in this part of the report draft:}
\begin{itemize}
    \item robots.txt support (\texttt{robots.py}),
    \item per-IP politeness/rate limiting (\texttt{politeness.py}),
    \item downloader (\texttt{downloader.py}),
    \item headless renderer (\texttt{renderer.py}),
    \item content-type classification and CLI test hooks (\texttt{main.py}, \texttt{config.py}).
\end{itemize}

\textbf{Implemented by teammate (integrated architecture context):}
\begin{itemize}
    \item DB-backed frontier and worker loop orchestration,
    \item parser for links/images and URL normalization,
    \item relevance scoring and preferential URL ordering,
    \item exact hash dedup integration in persistence layer.
\end{itemize}

\section{Preferential Crawling Strategy}
The crawler uses a preferential strategy, not simple BFS/DFS. Every discovered URL gets a relevance score, and workers always claim the highest-scored frontier item first.

Scoring is domain-specific. The current policy combines signals such as:
\begin{itemize}
    \item topic or section keywords in URL path,
    \item same-domain bonus for MedOverNet scope,
    \item penalties for low-value or off-topic paths,
    \item depth-related penalties for very deep paths,
    \item penalties for binary/document-heavy URLs when not relevant.
\end{itemize}

\textbf{Example policy behavior (illustrative):}
a forum or health-topic URL under MedOverNet gets higher priority than a generic file/document URL. Because frontier selection is ordered by score descending, workers crawl likely relevant pages first.

\textbf{PLACEHOLDER (final report):} add one concrete scored URL example from a real crawl run (input URL, feature contributions, final score).

\section{Robots.txt and Politeness}
The crawler enforces robots and ethics constraints before each fetch:
\begin{itemize}
    \item robots fetched per origin and cached,
    \item \texttt{Allow}/\texttt{Disallow} checked for crawler user-agent,
    \item \texttt{Crawl-delay} read when present,
    \item \texttt{Sitemap} URLs collected for optional frontier seeding,
    \item hard minimum delay of 5 seconds per target IP, even if robots does not define delay.
\end{itemize}

The implementation uses a thread-safe per-IP limiter. Effective delay is:
\[
    \text{effective\_delay} = \max(5s, \text{robots\_crawl\_delay})
\]

This design allows parallel workers while still respecting per-target pacing.

\textbf{PLACEHOLDER (final report):} insert one short robots compliance example from logs (allowed URL, disallowed URL, applied delay).

\section{Downloader, Renderer, and Content-Type Handling}
\subsection{HTTP downloader}
The downloader uses \texttt{requests} with crawler user-agent and timeout settings from configuration. It follows redirects and stores status code and final URL.

\subsection{JS rendering}
For pages where static HTML is not enough, the crawler can render using headless Chrome via Selenium. This mode is optional because it is slower, and we use it when needed for dynamic content.

\subsection{Content-type policy}
The crawler distinguishes HTML and binary resources:
\begin{itemize}
    \item HTML: stored in \texttt{page.html\_content}, \texttt{page\_type = HTML}
    \item PDF: supported as optional binary download, otherwise metadata-only
    \item DOC/DOCX/PPT/PPTX: metadata listed in \texttt{page\_data}, no binary payload required
    \item other non-HTML resources: marked as \texttt{BINARY}
\end{itemize}

This follows assignment constraints: crawl HTML content by default, handle PDF where needed, and avoid storing unnecessary binary blobs.

\section{Deduplication}
\subsection{Exact deduplication (implemented)}
The crawler computes a content hash (SHA-256 style exact hash) for HTML pages. If the hash already exists:
\begin{itemize}
    \item new page is marked \texttt{DUPLICATE},
    \item \texttt{duplicate\_of} points to the original page,
    \item duplicate page does not store full HTML content.
\end{itemize}

The schema includes \texttt{content\_hash} and \texttt{duplicate\_of} support for this flow.

\subsection{Partial deduplication (bonus, planned/not final yet)}
Locality-sensitive hashing + Jaccard similarity is planned as a bonus extension for near-duplicate detection. This is \textbf{not claimed as completed} in this draft.

\textbf{PLACEHOLDER (final report):} if implemented later, describe final LSH parameters and one near-duplicate example.

\section{Parameters and Configuration}
Crawler behavior is controlled by configuration variables. Main parameters include:
\begin{itemize}
    \item number of workers,
    \item crawler user-agent (\texttt{fri-wier-\{GROUP\_NAME\}}),
    \item minimum delay (must remain $\geq 5$ s),
    \item robots timeout,
    \item download timeout,
    \item rendering timeout,
    \item optional PDF download flag.
\end{itemize}

The command-line entry point in \texttt{main.py} was used to test:
\begin{itemize}
    \item robots checks and effective delay behavior,
    \item downloader behavior for HTML/PDF cases,
    \item renderer output size difference on dynamic pages.
\end{itemize}

\textbf{PLACEHOLDER (final report):} include final parameter table used for the 5000-page crawl run.

\section{Problems and Solutions}
\subsection{Robots and ethical crawling}
\textbf{Problem:} parallel workers can violate politeness if delays are not shared.\\
\textbf{Solution:} thread-safe per-IP limiter with effective delay max(5s, crawl-delay).

\subsection{Dynamic content}
\textbf{Problem:} some links/content are injected by JavaScript.\\
\textbf{Solution:} optional Selenium rendering path for selected pages.

\subsection{Frontier concurrency}
\textbf{Problem:} multiple workers may claim the same URL.\\
\textbf{Solution:} DB transaction with \texttt{FOR UPDATE SKIP LOCKED} and in-progress marking.

\subsection{Duplicate handling}
\textbf{Problem:} same content may appear under different URLs.\\
\textbf{Solution:} exact hash check and \texttt{DUPLICATE} page type with pointer to original.

\subsection{Spider traps and loops}
\textbf{Problem:} repeated navigation patterns and low-value loops can waste crawl budget.\\
\textbf{Solution:} canonical URL uniqueness in DB, relevance scoring, and robots compliance constraints.

\section{Results, Statistics, and Visualization (Placeholders)}
This section is intentionally a scaffold for the final submission stage.

\subsection{Crawldb statistics}
\textbf{PLACEHOLDER:} fill values from SQL queries after final crawl run.

\begin{table}[H]
    \centering
    \caption{Main crawl statistics (seed domain and full crawldb)}
    \begin{tabular}{@{}lll@{}}
        \toprule
        Metric                       & Seed domain (MedOverNet) & Full crawldb   \\ \midrule
        Number of sites              & \texttt{[TBD]}           & \texttt{[TBD]} \\
        Number of HTML pages         & \texttt{[TBD]}           & \texttt{[TBD]} \\
        Number of duplicates         & \texttt{[TBD]}           & \texttt{[TBD]} \\
        Number of binary pages       & \texttt{[TBD]}           & \texttt{[TBD]} \\
        Number of images             & \texttt{[TBD]}           & \texttt{[TBD]} \\
        Average images per HTML page & \texttt{[TBD]}           & \texttt{[TBD]} \\ \bottomrule
    \end{tabular}
\end{table}

\begin{table}[H]
    \centering
    \caption{Binary documents by type}
    \begin{tabular}{@{}ll@{}}
        \toprule
        Data type & Count          \\ \midrule
        PDF       & \texttt{[TBD]} \\
        DOC       & \texttt{[TBD]} \\
        DOCX      & \texttt{[TBD]} \\
        PPT       & \texttt{[TBD]} \\
        PPTX      & \texttt{[TBD]} \\ \bottomrule
    \end{tabular}
\end{table}

\subsection{Link graph visualization}
\textbf{PLACEHOLDER:} add graph visualization figure generated in Gephi/D3/vis.js.

\begin{figure}[H]
    \centering
    \fbox{\parbox{0.95\columnwidth}{\centering PLACEHOLDER: Link graph figure (nodes=pages, edges=links)}}
    \caption{Crawl link structure (placeholder).}
    \label{fig:link_graph_placeholder}
\end{figure}

\subsection{Example pages/images}
\textbf{PLACEHOLDER:} add 1--2 representative screenshots or extracted image examples.

\section{Conclusion}
We prepared a working foundation for a preferential MedOverNet crawler with core requirements: robots compliance, politeness control, HTML/binary handling, optional JS rendering, exact hash dedup support, and a multi-worker DB-backed crawl design.

The next step is the final full crawl run, followed by final statistics, link visualization, and optional near-duplicate (LSH + Jaccard) extension if completed. This report draft is ready for those final additions.

\end{document}
